\documentclass[aspectratio=169]{beamer}
\usepackage[utf8]{inputenc}
\usepackage[T1]{fontenc}
\usepackage[spanish]{babel}
\usepackage{tikz}
\usepackage{booktabs}
\usetikzlibrary{arrows.meta, positioning}

\usetheme{Madrid}
\usecolortheme{seahorse}

\title{Helix ALM}
\author{Jordy Fabi\'an Rivas Bodero}
\institute{UTEQ --- Ingenier\'ia de Requerimientos (ISR-401)}
\date{Noviembre 2026}

\begin{document}

\begin{frame}
  \titlepage
\end{frame}

\begin{frame}{Helix ALM: Evaluaci\'on de Adopci\'on}
  \begin{block}{Uso en Equipos \'Agiles}
    \begin{itemize}
      \item No conviene imponer Helix ALM como eje diario cuando se busca iteraci\'on r\'apida y baja carga administrativa \cite{pohl2015requirements}.
      \item Implica costo de licenciamiento, curva de aprendizaje y mantenimiento de flujos y enlaces \cite{pohl2015requirements}.
      \item En entornos livianos resulta preferible un flujo \'agil directo con \textbf{Git + Markdown + Issues}.
    \end{itemize}
  \end{block}
  \vspace{0.2cm}
  \begin{alertblock}{Entornos Cr\'iticos o Regulados}
    \begin{itemize}
      \item Resulta adecuado cuando se exigen auditor\'ias formales, cumplimiento normativo o trazabilidad estricta \cite{pohl2015requirements}.
    \end{itemize}
  \end{alertblock}
\end{frame}

\begin{frame}{Helix ALM: Veredicto T\'ecnico}
  \begin{columns}[T]
    \begin{column}{0.48\textwidth}
      \begin{block}{Cu\'ando Evitarlo}
        \begin{itemize}
          \item Equipos peque\~nos o prototipos con alta velocidad de cambio.
          \item Cuando la configuraci\'on de \'items supera el beneficio operativo.
          \item Presupuestos ajustados sin necesidad de certificaci\'on formal.
        \end{itemize}
      \end{block}
    \end{column}
    \begin{column}{0.48\textwidth}
      \begin{block}{Cu\'ando Aplicarlo}
        \begin{itemize}
          \item Gesti\'on centralizada de requisitos, casos de prueba e incidencias \cite{pohl2015requirements}.
          \item Trazabilidad bidireccional auditable entre artefactos \cite{pohl2015requirements}.
          \item An\'alisis de impacto ante modificaciones \cite{pohl2015requirements}.
        \end{itemize}
      \end{block}
    \end{column}
  \end{columns}
\end{frame}

\begin{frame}{Helix ALM: Trazabilidad de \'Items}
  \begin{center}
    \begin{tikzpicture}[
      node distance=1.4cm,
      box/.style={rectangle, draw=blue!60!black, fill=blue!8, thick, rounded corners, minimum width=3.4cm, minimum height=1.5cm, align=center, font=\small},
      arrow/.style={-{Stealth[length=3mm]}, thick, blue!70!black}
    ]
      \node[box] (req) {\textbf{Requisito (REQ-01)}\\Especificaci\'on de regla\\o condici\'on funcional\\del sistema};
      \node[box, right=of req] (tc) {\textbf{Caso de Prueba (TC-01)}\\Verificaci\'on del\\comportamiento y\\criterios de salida};
      \node[box, right=of tc] (inc) {\textbf{Incidencia (INC-01)}\\Registro de defecto o\\desviaci\'on detectada\\en ejecuci\'on};

      \draw[arrow] (req) -- node[above, font=\scriptsize] {Verifica} (tc);
      \draw[arrow] (tc) -- node[above, font=\scriptsize] {Genera} (inc);
      \draw[arrow, dashed, red!70!black] (inc.south) to[bend left=22] node[below, font=\scriptsize] {Impacta / Revisi\'on} (req.south);
    \end{tikzpicture}
  \end{center}

  \vspace{0.2cm}
  \begin{itemize}
    \item Enlaza requisitos, casos de prueba e incidencias dentro de un mismo repositorio de gesti\'on \cite{pohl2015requirements}.
  \end{itemize}
\end{frame}

\begin{frame}{Helix ALM: Matriz de Impacto}
  \begin{table}
    \centering
    \small
    \begin{tabular}{llll}
      \toprule
      \textbf{Artefacto} & \textbf{ID} & \textbf{Estado Anterior} & \textbf{Acci\'on por Alerta de Impacto} \\
      \midrule
      Requisito & REQ-01 & Aprobado (v1.0) & Actualizado a v1.1 (nuevo criterio) \\
      Caso de Prueba & TC-01 & Aprobado & Marcado como \textit{Sospechoso} (re-ejecutar) \\
      Incidencia & INC-01 & En revisi\'on & Reevaluar causa y alcance del cambio \\
      \bottomrule
    \end{tabular}
  \end{table}

  \vspace{0.2cm}
  \begin{block}{Propagaci\'on de Cambios}
    Helix ALM marca como \textit{sospechosos} los \'items vinculados cuando se modifica un elemento relacionado \cite{pohl2015requirements}.
  \end{block}
\end{frame}

\begin{frame}{Helix ALM: Conclusiones}
  \begin{block}{S\'intesis}
    \begin{itemize}
      \item Helix ALM aporta trazabilidad formal y control de impacto sobre requisitos, pruebas e incidencias \cite{pohl2015requirements}.
      \item Su adopci\'on debe evaluarse frente al costo operativo y al nivel de formalidad exigido \cite{pohl2015requirements}.
    \end{itemize}
  \end{block}
\end{frame}

\begin{frame}[allowframebreaks]{Referencias}
  \footnotesize
  \bibliographystyle{IEEEtran}
  \bibliography{referencias/referencias}
\end{frame}

\end{document}